\documentclass[12pt]{article}
\usepackage{amsmath, amssymb, array, booktabs, caption}
\usepackage{geometry}% 1-inch margins
\geometry{letterpaper, margin=1in}% 1-inch margins

\title{MATH 374 (Discrete Structures)\\Chapter 1.1 \& 1.2 Homework}
\author{Jacob Vaught}
\date{09-11-2023}

\begin{document}
\maketitle

\newpage
\section*{Problem 1.1.21}
Prove the following tautologies by starting with the left side and finding a series of equivalent wffs that will convert the left side into the right side. You may use any of the equivalences in the list on page 8 or in Exercise 20.
\begin{enumerate}
\item $(A \land B') \land C \leftrightarrow (A \land C) \land B'$
\item $(A \lor B) \land (A\lor B') \leftrightarrow A$
\item $A \land (A \land B')' \leftrightarrow A \land B$
\end{enumerate}


\newpage
\subsection*{Solutions}

\subsubsection*{1. Proving $(A \land B') \land C \leftrightarrow (A \land C) \land B'$}

\begin{align*}
(A \land B') \land C &= (A \land C) \land B' && \text{[Associative Law]} \\
\end{align*}
The statement is proven to be a tautology by direct application of the Associative Law.

\subsubsection*{2. Proving $(A \lor B) \land (A \lor B') \leftrightarrow A$}

\begin{align*}
(A \lor B) \land (A \lor B') &= A \lor (B \land B') && \text{[Distributive Law]} \\
&= A \lor \text{F} && \text{[Complement Law]} \\
&= A && \text{[Identity Law]} \\
\end{align*}
The statement is proven to be a tautology by sequentially applying the Distributive, Complement, and Identity Laws.

\subsubsection*{3. Proving $A \land (A \land B')' \leftrightarrow A \land B$}

\begin{align*}
A \land (A \land B')' &= A \land (A' \lor B) && \text{[DeMorgan's Law]} \\
&= A \land A' \lor A \land B && \text{[Distributive Law]} \\
&= \text{F} \lor A \land B && \text{[Complement Law]} \\
&= A \land B && \text{[Identity Law]} \\
\end{align*}
The left-hand side successfully converts into the right-hand side, proving the statement to be a tautology.


\newpage
\section*{Problem 1.1.31}
From the truth table for A $\lor$ B, the value of A $\lor$ B is true if A is true, if B is true, or if both are true. This use of the word "or" is the \textit{inclusive or}, where the result is true if both components are true, as in "We may have rain or drizzle tomorrow." Another use of the word "or" in the English language is the \textit{exclusive or}, sometimes written XOR, in which the result is false when both components are true. The exclusive or is understood in the sentence "At the intersection, you should turn north or south." Exclusive or is symbolized by A $\oplus$ B.
\begin{enumerate}
\item Write the truth table for the exclusive or.
\item Show that A $\oplus$ B $\leftrightarrow$ (A $\leftrightarrow$ B)' is a tautology.
\end{enumerate}


\newpage
\subsection*{Solution}

\begin{enumerate}
  \item* Write the truth table for the exclusive or.
  \section*{Solution}
  \begin{center}
    \begin{tabular}{|c|c|c|}
      \hline
      \( A \) & \( B \) & \( A \oplus B \) \\
      \hline
      0 & 0 & 0 \\
      0 & 1 & 1 \\
      1 & 0 & 1 \\
      1 & 1 & 0 \\
      \hline
    \end{tabular}
  \end{center}

  \item* Show that \( A \oplus B \leftrightarrow (A \leftrightarrow B)' \) is a tautology.
  \section*{Solution}
  \subsubsection*{Proof}
  To prove that \( A \oplus B \leftrightarrow (A \leftrightarrow B)' \) is a tautology, we can use a truth table.

  \begin{center}
    \begin{tabular}{|c|c|c|c|c|}
      \hline
      \( A \) & \( B \) & \( A \oplus B \) & \( A \leftrightarrow B \) & \( A \oplus B \leftrightarrow (A \leftrightarrow B)' \) \\
      \hline
      0 & 0 & 0 & 1 & 1 \\
      0 & 1 & 1 & 0 & 1 \\
      1 & 0 & 1 & 0 & 1 \\
      1 & 1 & 0 & 1 & 1 \\
      \hline
    \end{tabular}
  \end{center}

  As you can see, the final column is all 1's, which means that \( A \oplus B \leftrightarrow (A \leftrightarrow B)' \) is indeed a tautology.
\end{enumerate}




\newpage
\section*{Problem 1.2.9}
Justify each step in the proof sequence of $A \land (B \rightarrow C) \rightarrow (B \rightarrow (A \land C))$
\begin{enumerate}
\item $A$ \hfill
\item $B \rightarrow C$ \hfill
\item $B$ \hfill
\item $C$ \hfill
\item $A \land C$ \hfill
\end{enumerate}


\newpage
\subsection*{Solution}
\subsubsection*{Proof Justification}

We denote \( A' \) as the negation of \( A \) and \( B' \) as the negation of \( B \). \( \to \) represents implication, and \( \land \) is the logical "AND" operation.

\subsection*{Logical Proof*}

\begin{enumerate}
    \item \( A' \land (B \to A) \to B' \) \hfill Original Statement
    
    \item \( A' \land (\lnot B \lor A) \to B' \) \hfill Expand Implication
    
    \item \( (A' \land \lnot B) \lor (A' \land A) \to B' \) \hfill Distribute \( A' \) Over \( \lnot B \lor A \)
    
    \item \( (A' \land \lnot B) \lor \text{False} \to B' \) \hfill Contradiction Simplification
    
    \item \( A' \land \lnot B \to B' \) \hfill Simplify
\end{enumerate}

\subsection*{Proof by Contrapositive*}

\begin{enumerate}
    \item \( B \to \lnot (A' \land \lnot B) \) \hfill Contrapositive Statement
    
    \item \( B \) \hfill Assume \( B \) is True
    
    \item \( \lnot B \to \text{False} \) \hfill Negate \( \lnot B \)
    
    \item \( A' \land \lnot B \to \text{False} \) \hfill Examine \( A' \land \lnot B \)
    
    \item \( \lnot (A' \land \lnot B) \to \text{True} \) \hfill Check Contrapositive
\end{enumerate}

By showing the contrapositive \( B \to \lnot (A' \land \lnot B) \) is true, we thereby prove that the original statement \( A' \land \lnot B \to B' \) is also valid.



\newpage
\section*{Problem 1.2.11}
Justify each step in the proof sequence of $A' \land B \land [B \rightarrow (A \lor C)] \rightarrow C$
\begin{enumerate}
\item $A'$
\item $B$
\item $B \rightarrow (A \lor C)$
\item $A \lor C$ 
\item $(A')' \lor C$ 
\item $A' \rightarrow C$ 
\item $C$
\end{enumerate}


\newpage
\subsection*{Solution}
\subsubsection*{Proof Justification}

To justify each step in the proof sequence, let's consider the following:

\begin{enumerate}
    \item \( A' \) \hfill Assumed as part of the compound statement \( A' \land B \land [B \rightarrow (A \lor C)] \)
    
    \item \( B \) \hfill Assumed as part of the compound statement \( A' \land B \land [B \rightarrow (A \lor C)] \)
    
    \item \( B \rightarrow (A \lor C) \) \hfill Assumed as part of the compound statement \( A' \land B \land [B \rightarrow (A \lor C)] \)
    
    \item \( A \lor C \) \hfill Modus Ponens using steps 2 and 3
    
    \item \( (A')' \lor C \) \hfill Logical Equivalence; \( A' \) is logically equivalent to \( (A')' \)
    
    \item \( A' \rightarrow C \) \hfill Disjunctive Syllogism on step 4 using step 1;
    
    \item \( C \) \hfill Modus Ponens using steps 1 and 6
\end{enumerate}
We can confirm that the statement \( A' \land B \land [B \rightarrow (A \lor C)] \rightarrow C \) is proven to be true.



\newpage
\section*{Problem 1.2.12-21}
In Exercises 12-22, use propositional logic to prove that the argument is valid.
\begin{enumerate}
    \item[12.] \( A' \land (B \to A) \to B' \)
    \item[15.] \( A' \land (A \lor B) \to B \)
    \item[16.] \( [A \to (B \to C)] \land (A \lor D') \land B \to (D \to C) \)
    \item[17.] \( (A' \to B') \land B \land (A \to C) \to C \)
    \item[19.] \( [A \to (B \to C)] \to [B \to (A \to C)] \)
    \item[21.] \( (A \to C) \land (C \to B') \land B \to A' \)
\end{enumerate}


\newpage
\subsection*{Solution for \#12}
\subsubsection*{Logical Proof}
\begin{enumerate}
    \item \( A' \land (B \to A) \to B' \) \hfill Original Statement
    \item \( A' \land (B' \lor A) \to B' \) \hfill Expand Implication
    \item \( (A' \land B') \lor (A' \land A) \to B' \) \hfill Distribute \( A' \) Over \( B' \lor A \)
    \item \( (A' \land B') \lor \text{False} \to B' \) \hfill Contradiction Simplification
    \item \( A' \land B' \to B' \) \hfill Simplify
\end{enumerate}

\subsubsection*{Proof by Contrapositive}
\begin{enumerate}
    \item \( B \to (A' \land B')' \) \hfill Contrapositive Statement
    \item \( B \) \hfill Assume \( B \) is True
    \item \( B' \to \text{False} \) \hfill Negate \( B' \)
    \item \( A' \land B' \to \text{False} \) \hfill Examine \( A' \land B' \)
    \item \( (A' \land B')' \to \text{True} \) \hfill Check Contrapositive
\end{enumerate}
By showing the contrapositive \( B \to (A' \land B')' \) is true, we thereby prove that the original statement \( A' \land B' \to B' \) is also valid.

\newpage
\subsection*{Solution for \#15}
\subsubsection*{Logical Proof}
\begin{enumerate}
    \item \( A' \land (A \lor B) \to B \) \hfill Original Statement
    \item \( (A' \land A) \lor (A' \land B) \to B \) \hfill Distribute \( A' \) over \( A \lor B \)
    \item \( \text{False} \lor (A' \land B) \to B \) \hfill \( A' \land A \) is a Contradiction, so it simplifies to \text{False}
    \item \( A' \land B \to B \) \hfill Elimination of the \text{False} term
\end{enumerate}

\subsubsection*{Proof by Contrapositive}
\begin{enumerate}
    \item \( B' \to (A' \land (A \lor B))' \) \hfill Contrapositive Statement
    \item \( B' \) \hfill Assume \( B' \) is True
    \item \( B' \to \text{False} \) \hfill Negate \( B' \)
    \item \( A \lor A' \) \hfill Examine \( A' \land (A \lor B) \)
    \item \( A' \lor A \to \text{True} \) \hfill Check Contrapositive
\end{enumerate}
By following the steps of both the Logical Proof and the Proof by Contrapositive, we have validated that the original statement \( A' \land (A \lor B) \to B \) is indeed valid.

% Problem Statement for #16
\newpage
\subsection*{Solution for \#16}

\subsubsection*{Logical Proof}

\begin{enumerate}
    \item \( [A \to (B \to C)] \land (A \lor D') \land B \to (D \to C) \) \hfill Original Statement
    \item \( (B' \lor C) \land (A \lor D') \land B \to (D' \lor C) \) \hfill Expand Implication
    \item \( ((B' \land A) \lor (B' \land D') \lor C) \to (D' \lor C) \) \hfill Distribute over \( B' \lor C \)
    \item \( (B' \land A) \lor (B' \land D') \lor C \to C \) \hfill Eliminate \( D' \lor C \) 
    \item \( C \to C \) \hfill Simplify
\end{enumerate}

\subsubsection*{Proof by Contrapositive}

\begin{enumerate}
    \item \( (D' \lor C') \to [(B' \lor C) \land (A \lor D')]' \) \hfill Contrapositive Statement
    \item \( D' \lor C' \) \hfill Assume \( D' \lor C' \) is True
    \item \( D \land C \to \text{False} \) \hfill Negate \( D' \lor C' \)
    \item \( (B' \lor C) \land (A \lor D') \to \text{False} \) \hfill Examine Original Statement
    \item \( [(B' \lor C) \land (A \lor D')]' \to \text{True} \) \hfill Check Contrapositive
\end{enumerate}

% Problem Statement for #17
\newpage
\subsection*{Solution for \#17}

\subsubsection*{Logical Proof}

\begin{enumerate}
    \item \( (A' \to B') \land B \land (A \to C) \to C \) \hfill Original Statement
    \item \( (A \lor B') \land B \land (A' \lor C) \to C \) \hfill Expand Implication
    \item \( ((A \lor B') \land B) \land (A' \lor C) \to C \) \hfill Distribute over \( A' \lor C \)
    \item \( ((A \lor B) \land B) \land (A' \lor C) \to C \) \hfill Simplify \( A \lor B' \) as \( A \lor B \)
    \item \( C \to C \) \hfill Simplify
\end{enumerate}

\subsubsection*{Proof by Contrapositive}

\begin{enumerate}
    \item \( C' \to [(A' \to B') \land B \land (A \to C)]' \) \hfill Contrapositive Statement
    \item \( C' \) \hfill Assume \( C' \) is True
    \item \( C \to \text{False} \) \hfill Negate \( C' \)
    \item \( (A' \to B') \land B \land (A \to C) \to \text{False} \) \hfill Examine Original Statement
    \item \( [(A' \to B') \land B \land (A \to C)]' \to \text{True} \) \hfill Check Contrapositive
\end{enumerate}

% Problem Statement for #19
\newpage
\subsection*{Solution for \#19}

\subsubsection*{Logical Proof}

\begin{enumerate}
    \item \( [A \to (B \to C)] \to [B \to (A \to C)] \) \hfill Original Statement
    \item \( (A' \lor (B' \lor C)) \to (B' \lor (A' \lor C)) \) \hfill Expand Implication
    \item \( A' \lor B' \lor C \to B' \lor A' \lor C \) \hfill Simplify
    \item \( B' \lor A' \lor C \to B' \lor A' \lor C \) \hfill Commutativity
\end{enumerate}

\subsubsection*{Proof by Contrapositive}

\begin{enumerate}
    \item \( (B' \lor (A' \lor C))' \to (A' \lor (B' \lor C))' \) \hfill Contrapositive Statement
    \item \( B \land (A \land C') \) \hfill Assume \( (B' \lor (A' \lor C))' \) is True
    \item \( B \land A \land C' \to \text{False} \) \hfill Negate \( B' \lor (A' \lor C) \)
    \item \( A' \lor (B' \lor C) \to \text{False} \) \hfill Examine Original Statement
    \item \( (A' \lor (B' \lor C))' \to \text{True} \) \hfill Check Contrapositive
\end{enumerate}

% Problem Statement for #21
\newpage
\subsection*{Solution for \#21}

\subsubsection*{Logical Proof}

\begin{enumerate}
    \item \( (A \to C) \land (C \to B') \land B \to A' \) \hfill Original Statement
    \item \( (A' \lor C) \land (C' \lor B') \land B \to A' \) \hfill Expand Implication
    \item \( ((A' \lor C) \land B) \land (C' \lor B') \to A' \) \hfill Distribute over \( A' \lor C \)
    \item \( (A' \lor C) \land B \land C' \lor B' \to A' \) \hfill Simplify
    \item \( A' \land B \land C' \lor B' \to A' \) \hfill Distribute \( B \) over \( C' \lor B' \)
    \item \( A' \to A' \) \hfill Simplify
\end{enumerate}

\subsubsection*{Proof by Contrapositive}

\begin{enumerate}
    \item \( A \to [(A' \lor C) \land (C' \lor B') \land B]' \) \hfill Contrapositive Statement
    \item \( A \) \hfill Assume \( A \) is True
    \item \( A' \to \text{False} \) \hfill Negate \( A' \)
    \item \( (A' \lor C) \land (C' \lor B') \land B \to \text{False} \) \hfill Examine Original Statement
    \item \( [(A' \lor C) \land (C' \lor B') \land B]' \to \text{True} \) \hfill Check Contrapositive
\end{enumerate}

\newpage
\section*{Problem 1.2.40}
Using propositional logic, including the rules in Table 1.14, prove that each argument in Exercises 40-48 is valid. Use the statement letters shown.
\newline

40. If the program is efficient, it executes quickly. Either the program is efficient, or it has a bug. However, the program does not execute quickly. Therefore it has a bug. E, Q, B

\newpage
\subsection*{Solution}
\subsubsection*{Derivation of Original Assumptions}
The problem provides three main premises:
\begin{enumerate}
    \item "If the program is efficient, it executes quickly" \hfill Translated as \( E \to Q \)
    \item "Either the program is efficient, or it has a bug" \hfill Translated as \( E \lor B \)
    \item "The program does not execute quickly" \hfill Translated as \( Q' \)
    \item "Therefore it has a bug." \hfill Translated as \( \to B \)
\end{enumerate}
We combine these premises to formulate the original assumptions: \( (E \to Q) \land (E \lor B) \land Q' \to B\).

\subsubsection*{Logical Proof}
\begin{enumerate}
    \item \( (E \to Q) \land (E \lor B) \land Q' \) \hfill Original Assumptions (Derived from problem statement)
    \item \( E \to Q \) \hfill Extract from Original Assumptions
    \item \( E \lor B \) \hfill Extract from Original Assumptions
    \item \( Q' \) \hfill Extract from Original Assumptions
    \item \( E' \lor Q \) \hfill Rewrite \( E \to Q \) in terms of \( \lor \)
    \item \( E' \lor Q' \) \hfill Apply \( Q' \) from the original assumptions
    \item \( E' \) \hfill Simplify \( E' \lor Q' \)
    \item \( B \) \hfill Since \( E' \), from \( E \lor B \), it must be \( B \)
\end{enumerate}

\subsubsection*{Conclusion}
Therefore, it has a bug (\( B \)).
\newline \newline
Through this logical proof, we have demonstrated that the original assumptions \( (E \to Q) \land (E \lor B) \land Q' \) lead to the conclusion that the program has a bug \( B \), which is in line with the argument "Therefore, it has a bug."


\newpage
\section*{Appendix}
\begin{table}[h]
    \centering
    \caption*{Table: Algebraic Properties of Logic Operators}
    \begin{tabular}{l l l }
        \toprule
        Property & Expression & Expression \\
        \midrule
        Commutative & \( A \land B \leftrightarrow B \land A \) &  \( A \lor B \leftrightarrow B \lor A \) \\
        Associative & \( (A \land B) \land C \leftrightarrow A \land (B \land C) \) & \( (A \lor B) \lor C \leftrightarrow A \lor (B \lor C) \) \\
        Distributive & \( A \land (B \lor C) \leftrightarrow (A \land B) \lor (A \land C) \) & \( A \lor (B \land C) \leftrightarrow (A \lor B) \land (A \lor C) \) \\
        Identity & \( A \land 0 \leftrightarrow A \) &  \( A \lor 1 \leftrightarrow A \) \\
        Complement & \( A \land A' \leftrightarrow 0 \) &  \( A \lor A' \leftrightarrow 1 \) \\
        \bottomrule
    \end{tabular}
\end{table}

\begin{table}[h]
    \centering
    \caption*{Table 1.13: Inference Rules}
    \begin{tabular}{l l l}
        \toprule
        From & Can Derive & Name/Abbreviation for Rule \\
        \midrule
        \(P, P \rightarrow Q\) & \(Q\) & Modus Ponens - mp \\
        \(P \rightarrow Q, Q'\) & \(P'\) & Modus Tollens - mt \\
        \(P, Q\) & \(P \land Q\) & Conjunction - con \\
        \(P \land Q\) & \(P, Q\) & Simplification - sim \\
        \(P\) & \(P \lor Q\) & Addition - add \\
        \bottomrule
    \end{tabular}
\end{table}

\begin{table}[h]
    \centering
    \caption*{Table 1.14: More Inference Rules}
    \begin{tabular}{l l l}
        \toprule
        From & Can Derive & Name/Abbreviation for Rule \\
        \midrule
        \(P \rightarrow Q, Q \rightarrow R\) & \(P \rightarrow R\) [Ex. 16] & Hypothetical syllogism - hs \\
        \(P \lor Q, P'\) & \(Q\) [Ex. 23] & Disjunctive syllogism - ds \\
        \(P \rightarrow Q\) & \(Q' \rightarrow P'\) [Ex. 24] & Contraposition - cont \\
        \(Q' \rightarrow P'\) & \(P \rightarrow Q\) [Ex. 25] & Contraposition - cont \\
        \(P\) & \(P \land P\) [Ex. 26] & Self-reference - self \\
        \(P \lor P\) & \(P\) [Ex. 27] & Self-reference - self \\
        \((P \land Q) \rightarrow R\) & \(P \rightarrow (Q \rightarrow R)\) [Ex. 28] & Exportation - exp \\
        \(P, P'\) & \(Q\) [Ex. 29] & Inconsistency - inc \\
        \(P \land (Q \lor R)\) & \((P \land Q) \lor (P \land R)\) [Ex. 30] & Distributive - dist \\
        \(P \lor (Q \land R)\) & \((P \lor Q) \land (P \lor R)\) [Ex. 31] & Distributive - dist \\
        \bottomrule
    \end{tabular}
\end{table}

\end{document}
